\chapter{Managing Resources}
\label{ch:gm-resources}
\index{resources!GM guidance}

\begin{tcolorbox}[colback=blue!5!white, colframe=blue!75!black, title={Philosophy of Resources}]
\index{resources!philosophy}
\textbf{Resources are not there to hoard -- they are tools to create drama.}

\begin{itemize}
  \item \textbf{Spend Boons freely.} Hoarding them for XP conversion starves the table of tension and robs players of their safety net.
  \item \textbf{Let Fatigue build.} Fatigue is not punishment; it is the rising tide that forces hard choices and heroic comebacks.
  \item \textbf{Obligation is plot, not punishment.} When a Patron calls in a debt, a story hook arrives. Embrace it.
\end{itemize}

The most memorable sessions come from characters spending their last Boon on a desperate re-roll, staggering forward under three Fatigue, or accepting a Patron's demand that changes everything.
\end{tcolorbox}

In \textbf{Fate's Edge}, resources are narrative levers, not inventory lists. From the last sip of water in sun-scorched wastes to the fragile loyalty of a Ykrul war-band, every resource tells a story, and every story demands its price. As GM, you are the keeper of these threads---the weaver of scarcity and abundance. This chapter shows how to manage the systems that fuel mortal ambitions and epic campaigns.

\section{Supply Clock: The Pulse of Survival}
\label{sec:supply-timer}
\index{Supply Clock}

The \textbf{Supply Clock} tracks the party's access to essentials: food, water, gear, and logistical support. It is a shared party resource---a narrative lever that tightens tension when characters are isolated, pursued, or cut off from civilization.

\subsection{Supply Timer States}
\label{subsec:supply-stats}

\begin{table}[htbp]
\centering
\small
\begin{tabularx}{\linewidth}{lX}
\toprule
\textbf{Segments Filled} & \textbf{Narrative Effects} \\
\midrule
0 (Full) & Well-equipped and prepared; no penalty. \\
1 (Low) & Minor complications: bland rations, damaged arrows, thinning waterskins. \\
2 (Dangerous) & Each character gains \textbf{Fatigue 1} as scarcity takes its toll. \\
3 (Empty) & Severe penalties: starvation, dehydration, equipment failure imminent. \\
\bottomrule
\end{tabularx}
\caption{Supply Timer conditions and their narrative effects}
\label{tab:supply-timer-states}
\end{table}

\subsection{When to Advance the Supply Clock}
\label{subsec:supply-advance}

The Supply Timer advances when the world turns against the party's preparations:
\begin{itemize}
  \item Extended travel through hostile lands without proper provisioning (see \textbf{Travel Framework}, Chapter 9 of the Player's Guide or SRD Section 41).
  \item The GM spends \textbf{2+ Story Beats} on logistics failures or environmental hardships.
  \item The party chooses to travel light for speed or stealth.
  \item Failed \textbf{Survival} or \textbf{Craft} rolls related to hunting, foraging, or repair.
\end{itemize}

\subsection{Clearing Supply Segments}
\label{subsec:supply-clear}

The timer recedes when the party finds respite:
\begin{itemize}
  \item Reaching civilization resets the timer to \textit{Full}.
  \item A group \textbf{Survival} check (Wits + Survival, DV 2) in favorable conditions clears 1 segment.
  \item Downtime spent in relative safety removes 1 segment.
  \item Successful provisioning (a good hunt, a discovered cache) can reduce segments.
\end{itemize}

\paragraph{Example.}
A week-long sea passage with uncertain winds. A failed \textbf{Navigation} roll causes the GM to spend 2 SB -- filling two segments as supplies spoil. The party is now at \textit{Low Supply}. A second failed roll fills another segment -- \textit{Dangerously Low}. Fatigue sets in. The sea gnaws at their endurance.

\section{Fatigue: The Weight of the World}
\label{sec:fatigue-gm}
\index{Fatigue!GM guidance}

Fatigue represents the cumulative toll of journeying through a world that rarely offers comfort: exhaustion that seeps into bones, hunger that hollows cheeks, strain that clouds judgment.

\begin{tcolorbox}[colback=gray!5!white, colframe=gray!75!black, title={The Roller-Coaster Effect}, label={sec:roller-coaster}]\index{Fatigue!roller-coaster effect}

\textbf{Taking Harm clears Fatigue temporarily} -- the shock of a wound can jolt a character back into action. This creates a natural \textquotedbl roller coaster\textquotedbl\ of tension:

\medskip
\noindent\textit{Example:} Kael is at Fatigue 3 (Desperate, --2 dice). A goblin stabs him for Harm 1. The Harm conversion procedure: first convert Harm $\rightarrow$ Fatigue (none, light armor), then apply remaining Harm -- which clears all Fatigue. Kael's Fatigue drops to 0, but his Harm increases to 1. He is wounded but suddenly clear-headed, ready to fight back.

\medskip
Use this effect to create dramatic swings: a character near collapse who takes a wound and finds a second wind. It keeps combat unpredictable and rewards players who embrace risk.
\end{tcolorbox}

\subsection{Fatigue Effects}
\label{subsec:fatigue-effects}

Each point of Fatigue worsens the character's \textbf{Position} by one step (Dominant $\rightarrow$ Controlled $\rightarrow$ Desperate). If already Desperate, each additional Fatigue instead imposes a \textbf{--1 die penalty} on all rolls. Penalties from Fatigue and Harm \textbf{stack}.

\begin{table}[htbp]
\centering
\small
\begin{tabularx}{\linewidth}{lX}
\toprule
\textbf{Fatigue Level} & \textbf{Effect} \\
\midrule
1 & Position worsens one step (Dominant $\rightarrow$ Controlled $\rightarrow$ Desperate) \\
2 & If already Desperate, --1 die; otherwise Position worsens \\
3 & --2 dice (if already Desperate from Fatigue) \\
4+ & Collapse, KO, or spiritual break (incapacitated) \\
\bottomrule
\end{tabularx}
\caption{Fatigue progression and associated mechanical effects}
\label{tab:fatigue-effects}
\end{table}

\subsection{Clearing Fatigue}
\label{subsec:fatigue-clear}

Recovery requires genuine respite:
\begin{itemize}
  \item A \textbf{short rest} (quiet watch, food/water) removes \textbf{2 Fatigue}.
  \item A \textbf{full night's rest} in safety removes \textbf{all Fatigue}.
  \item Fatigue cannot be removed while the party is at \textit{Dangerous} or \textit{Empty} Supply.
  \item Medical attention (\textbf{Presence + Medicine}, DV 2) during downtime can remove an additional 1 Fatigue.
\end{itemize}

\paragraph{Narrative note:}
Fatigue is not only physical -- it can reflect mental strain, grief, or spiritual exhaustion. A failed ritual might leave a caster at Fatigue 2 from metaphysical backlash alone.

\subsection{Fatigue $\rightarrow$ Harm Conversion}
\label{subsec:fatigue-harm}

Each character has a Fatigue track equal to their \textbf{Body} attribute. When the track fills:
\begin{enumerate}
  \item Increase \textbf{Harm} by one level (0$\rightarrow$Harm 1, Harm 1$\rightarrow$Harm 2, Harm 2$\rightarrow$Harm 3).
  \item Clear all Fatigue.
\end{enumerate}
This conversion can occur multiple times in a scene. (See the \textit{Player's Guide} for full Harm rules.)

\section{Boons: The Currency of Resilience}
\label{sec:boons-gm}
\index{Boons!GM guidance}

Boons are narrative tokens earned by embracing risk. They reward \textit{failure with opportunity} -- the silver lining in clouds of defeat.

\begin{tcolorbox}[colback=yellow!5!white, colframe=yellow!75!black, title={Why You Should Spend Boons -- The Math of Fun}]
\index{Boons!spending encouragement}
\textbf{Hoarding Boons for XP conversion leads to less engagement and lower fun scores.}

\medskip
\noindent Every Boon you hold is a dramatic moment waiting to happen:
\begin{itemize}
  \item Re-rolling a critical failure turns despair into triumph.
  \item Improving Position before a risky roll turns a Desperate gamble into a Controlled chance.
  \item Activating an Asset brings your character's network into the story.
\end{itemize}
\textbf{XP from Boons is a trickle (max 2 per session).} The real reward is \textit{agency}. The player who spends Boons shapes the story; the player who hoards them watches from the sidelines.

\medskip
\noindent As GM, remind players: \textquotedbl You have 3 Boons. What will you spend them on?\textquotedbl\ When a player sits on 5 Boons, gently suggest they use one before the scene ends. A spent Boon is a moment of player power. An unspent Boon is wasted potential.
\end{tcolorbox}

\subsection{Earning Boons (GM Reminders)}
\label{subsec:boons-earn}

\begin{itemize}
  \item On a \textbf{Partial} (1 Boon) or \textbf{Miss} (2 Boons) with meaningful consequences.
  \item Through clever or risky roleplay that drives the story.
  \item Via bond-driven actions with intricate descriptions.
  \item At GM discretion for exceptional collaborative play.
\end{itemize}
\textbf{Limit:} At most \textbf{2 Boons} from failures per character per scene. Repetition of the same approach with the same stakes does not award additional Boons.

\subsection{Boon Economy (Quick Reference)}
\label{subsec:boons-economy}

\begin{itemize}
  \item \textbf{Holding cap:} Maximum 5 Boons.
  \item \textbf{Carryover:} At the end of each scene, reduce held Boons to a maximum of 2 (excess lost).
  \item \textbf{Conversion:} Once per session during downtime, convert 2 Boons $\rightarrow$ 1 XP (max 2 XP per session via conversion).
\end{itemize}

\subsection{Spending Boons (Player Options)}
\label{subsec:boons-spend}

Players may spend 1 Boon to:
\begin{itemize}
  \item Re-roll a single die (after seeing the pool).
  \item Activate an on-screen Asset.
  \item Improve Position by one step before a roll.
  \item Power certain Rites or abilities as described in their text.
\end{itemize}

\paragraph{Substituting Fatigue for Boons (Optional)}
If a player has no Boons remaining, they may substitute \textbf{1 Fatigue} for 1 Boon \textit{once per scene}. This represents pushing through exhaustion for a last chance. The GM may veto if the fiction does not support it.

\subsection{Anti-Fishing Measures (GM Guidance)}
\label{subsec:boons-anti-fishing}

\begin{itemize}
  \item \textbf{Once/Scene Cap:} At most 2 Boons from failures per character per scene.
  \item \textbf{Repetition Rule:} Same approach + same stakes in the same scene cannot award another Boon.
  \item \textbf{Position Gate:} \textit{Dominant} tests with trivial fallout do not award Boons.
\end{itemize}
\emph{Design note:} Boons reward narrative risk, not repeated failure.

\section{Followers and Assets: Power Beyond the Self}
\label{sec:followers-assets-gm}
\index{Followers!GM guidance}\index{Assets!GM guidance}

Players can invest XP in \textbf{Followers} (on-screen allies) and \textbf{Assets} (off-screen resources). These are story agents with their own motivations and risks. The following sections present simplified, low-bookkeeping systems that avoid numeric timers while preserving meaningful consequences.

\subsection{Followers: On-Screen Allies}
\label{subsec:followers-gm}

Followers are bought with XP and tracked by a \textbf{Cap} (maximum assist bonus). Cost = Cap\(^2\) XP; downtime 1--3 days to recruit.

\paragraph{Follower States}
Instead of numeric timers, track a follower's \textbf{Loyalty} and \textbf{Fitness} as one of three states each.

\begin{table}[htbp]
\centering
\small
\begin{tabularx}{\linewidth}{lX}
\toprule
\textbf{State} & \textbf{Mechanical Effect} \\
\midrule
\multicolumn{2}{c}{\textit{Loyalty}} \\
Faithful & Normal assist, independent actions available \\
Strained & Assist at –1 die; may refuse risky orders \\
Broken & Follower leaves or becomes hostile \\
\midrule
\multicolumn{2}{c}{\textit{Fitness}} \\
Ready & Full function \\
Hurt & Assist at –1 die; no independent actions \\
Down & Incapacitated; requires rescue or medical care \\
\bottomrule
\end{tabularx}
\caption{Follower Loyalty and Fitness states}
\label{tab:follower-states}
\end{table}

Change states when the fiction demands it (e.g., after a failed mission, a moral compromise, or taking damage). No rolls required.

\paragraph{Assisting in Scenes}
\begin{itemize}
  \item Assist dice equal \(\min(\text{Cap}, \text{helper's relevant Skill})\).
  \item Total assist from all sources capped at \textbf{+3 dice} (the \textit{Exceptional Coordination} talent allows +4 from one follower).
  \item Only one follower may assist a given action.
\end{itemize}

\paragraph{Follower Independent Actions}
Once per scene, \textbf{one} follower may take a small independent action (party-wide limit). Examples:
\begin{itemize}
  \item \textbf{Scout \& Signal} -- Change an ally's next action to \textit{Dominant}.
  \item \textbf{Distract \& Draw} -- Reduce a threat/hunt/escape timer by 1 tick.
  \item \textbf{Fetch \& Carry} -- Move a small object through danger.
\end{itemize}
\textbf{Cost:} Mark \textbf{Exposure +1} or worsen \textbf{Fitness} one step (Ready $\rightarrow$ Hurt $\rightarrow$ Down) on that follower.

\paragraph{Followers as Successors}
A Follower (Cap 3+) can become a \textbf{successor} if the player's main character retires, dies, or ascends. This is the core of the Legacy Engine (see \ref{ch:legacy-engine} in the full GM Guide).

\paragraph{Follower Upkeep and Risk}
\begin{itemize}
  \item Each downtime, pay XP equal to Cap \textbf{or} spend a scene tending the relationship.
  \item Neglect $\rightarrow$ Loyalty worsens one step (Faithful $\rightarrow$ Strained $\rightarrow$ Broken).
  \item If the GM spends \textbf{2+ SB} on an action you take with assistance, they may worsen Loyalty or Fitness on the follower instead of other consequences.
  \item \textbf{Off-Screen Capability:} Once per downtime, a Cap 3+ follower can solve a significant problem but generates \textbf{1 SB} for the party (the GM describes the resulting complication).
\end{itemize}

\subsection{Assets: Off-Screen Influence}
\label{subsec:assets-gm}

Assets are titles, safehouses, spy rings, charters. They do not act in scenes directly but change the fiction between sessions.

\begin{table}[htbp]
\centering
\small
\begin{tabularx}{\linewidth}{lX}
\toprule
\textbf{Tier} & \textbf{Cost and Examples} \\
\midrule
Minor & 4 XP (1 day) -- safehouse, small shop, local contact \\
Standard & 8 XP (1 week) -- noble title, spy ring, workshop \\
Major & 12 XP (1 month) -- city license, fortress lease, regional network \\
\bottomrule
\end{tabularx}
\caption{Asset tiers, acquisition cost, and representative examples}
\label{tab:asset-tiers}
\end{table}

\paragraph{Asset Status}
Assets have a single unified \textbf{Status} that describes their overall health and usability.

\begin{table}[htbp]
\centering
\small
\begin{tabularx}{\linewidth}{lX}
\toprule
\textbf{Status} & \textbf{Effect} \\
\midrule
Flourishing & +1 die on rolls involving the asset, +1 Yield \\
Stable & No modifier \\
Strained & –1 die on asset-related rolls, no Yield \\
Collapsed & Non-functional; requires quest or major downtime to restore \\
\bottomrule
\end{tabularx}
\caption{Asset Status levels and their effects}
\label{tab:asset-status}
\end{table}

The Status worsens when the asset is neglected, attacked, or suffers a public failure. It improves through dedicated downtime actions or successful missions that benefit the asset.

\paragraph{Using Assets}
\begin{itemize}
  \item \textbf{Off-Screen Effect:} Use once per session for free (the quiet work between adventures).
  \item \textbf{On-Screen Activation:} Spend 1 Boon to use the asset dramatically in a scene.
  \item \textbf{Downtime Activation:} Pay 2 XP or 1 Boon to activate an asset at campaign start or during downtime.
  \item The asset must have plausible scope for the intended effect.
\end{itemize}

\paragraph{Asset Upkeep}
Each downtime:
\begin{itemize}
  \item \textbf{Efficient (higher XP, less time):} Cost = \(\max(1, \frac{\text{Acquisition XP}}{3})\). Minimal effort.
  \item \textbf{Intensive (lower XP, more time):} Cost = 1 XP. Requires a dedicated downtime action with significant personal involvement.
  \item Failure to pay $\rightarrow$ Status worsens one step (Flourishing $\rightarrow$ Stable $\rightarrow$ Strained $\rightarrow$ Collapsed).
\end{itemize}

\section{Terrestrial Patrons: Favors and Debts}
\label{sec:terrestrial-patrons-gm}
\index{Patrons!terrestrial}\index{Debt Clock}

Not every patron is a cosmic entity. Powerful mortals---nobles, crime lords, generals, merchant princes---can act as \textbf{terrestrial patrons}. They offer tangible aid (escorts, legal protection, supplies, even ongoing Assets) in exchange for favors.

\subsection{How Terrestrial Patrons Work}
When a character accepts a favor from a terrestrial patron, start a \textbf{Debt Clock} with the patron's name (e.g., \textquotedbl Prefect Marcellus's Favor [4]\textquotedbl). Mark one segment each time the character accepts a significant favor. When the timer fills, the patron calls in the debt -- a mission, a political favor, or a direct demand.

The party also tracks a separate \textbf{Obligation} tally for each terrestrial patron, exactly as they would for a supernatural Patron (see \ref{sec:obligation-corruption-gm}). Each accepted favor adds +1 Obligation. Exceeding Obligation Capacity (Spirit + Presence) causes Fatigue exactly as with cosmic Patrons.

Refusing to honor the debt turns the patron into an enemy. The GM may advance associated \textit{Crisis} timers or introduce complications.

\begin{tcolorbox}[colback=gray!5!white, colframe=gray!75!black, title={Use Debt Timers Aggressively}]
\index{Debt Clock!GM advice}
\textbf{Debt timers are story fuel.} Call them in at the worst possible moment:
\begin{itemize}
  \item \textbf{Early (1--2 segments):} A messenger arrives with a \textquotedbl reminder.\textquotedbl\ Small favor.
  \item \textbf{Mid (3--4 segments):} A significant service is demanded.
  \item \textbf{Full (5--6 segments):} The patron calls in the debt at a critical moment -- refusal means enmity.
\end{itemize}
Debts that are never called are wasted potential.
\end{tcolorbox}

\subsection{Clearing Terrestrial Obligation}
Obligation to a terrestrial patron is cleared by performing \textbf{tasks} for the patron. The GM sets the scope:
\begin{itemize}
  \item \textbf{Minor task} (deliver a message, pay a bribe): clears 1 Obligation.
  \item \textbf{Significant task} (escort a caravan, break a siege): clears 2 Obligation.
  \item \textbf{Major task} (assassinate a rival, broker a treaty): clears 3+ Obligation.
\end{itemize}
Each task becomes a story hook.

\subsection{Examples of Terrestrial Patron Debts}
\begin{tcolorbox}[title=Sample Terrestrial Patrons, colback=white!97!gray, colframe=black!80!gray]
\begin{tabularx}{\linewidth}{lX}
\toprule
\textbf{Patron} & \textbf{Typical Favors and Debt Trigger} \\
\midrule
Prefect Marcellus (Ecktorian commander) & Safe passage, food, guards. Debt: dangerous mission behind enemy lines. \\
Censor Cassia (judicial authority) & Legal immunity, archive access. Debt: expose a political rival. \\
Khatun Sarnai (Ykrul chieftain) & Remounts, safe camp, guides. Debt: blood-price or raid on a rival clan. \\
Mistress Ilvar (Silkstrand crime lord) & Black Bazaar access, cutpurses. Debt: public assassination. \\
\bottomrule
\end{tabularx}
\end{tcolorbox}

\subsection{Running Terrestrial Patrons}
\begin{itemize}
  \item Track each patron's Debt Timer and the party's Obligation visibly.
  \item When a player asks for a favor, state the cost: \textquotedbl That's worth one segment and +1 Obligation. Are you willing?\textquotedbl
  \item When a Debt Timer fills, give a hook that fits the patron's personality.
  \item If the party refuses, tick \textit{Crisis} or introduce a complication.
\end{itemize}

\section{Player Trusts: Shared Legacy}
\label{sec:trusts-gm}
\index{Trusts}\index{shared XP pool}

A \textbf{Trust} is a shared pool of XP that players contribute to, used to acquire and maintain \textbf{party-owned Assets and Followers}. Trusts represent a collective enterprise---a mercenary company, a trading guild, a spy network, a religious order.

\subsection{Forming a Trust}
Any player may contribute XP to the Trust pool. Contributions are voluntary and cannot be withdrawn. The Trust has no individual owner; spending from the pool requires majority agreement.

\subsection{What a Trust Can Buy}
\begin{table}[htbp]
\centering
\small
\begin{tabularx}{\linewidth}{lX}
\toprule
\textbf{Cost} & \textbf{Benefit} \\
\midrule
4 XP & Minor Asset (shared safehouse, workshop, contact network) \\
8 XP & Standard Asset (guild charter, spy ring, fortress lease) \\
9 XP & Follower (Cap 3, loyal to the Trust, not any single PC) \\
\bottomrule
\end{tabularx}
\caption{Trust expenditures}
\label{tab:trust-costs}
\end{table}

Assets and Followers acquired through a Trust belong to the \textbf{party as a whole}. Any member may activate them (spending their own Boons).

\subsection{Upkeep and Delegation}
Each downtime, the Trust must pay upkeep for its Assets and Followers using the standard rules (1 XP or a scene per asset/follower). Alternatively, the Trust may \textbf{delegate upkeep} to a \textbf{Cap 5 Follower}:
\begin{itemize}
  \item The Cap 5 follower can manage up to \textbf{4 Assets} (or any combination of Assets and Followers totaling 4 slots).
  \item The follower's own upkeep covers the delegated Assets---no separate XP cost.
  \item \textbf{Risk:} If the follower becomes \textit{Strained} or \textit{Broken} (Loyalty), all delegated Assets worsen one step or are lost until the follower is restored.
\end{itemize}

\subsection{Dissolving a Trust}
If the party disbands, remaining XP in the pool is lost. Assets may be transferred to individual characters who pay the XP cost (or accept Obligation to a Terrestrial Patron to cover the difference).

\section{Campaign Resources: Mandate and Crisis}
\label{sec:campaign-resources}
\index{Mandate}\index{Crisis}

Two campaign-level timers track the party's influence and the world's resistance.

\subsection{Mandate Timer (0--6)}
\label{subsec:mandate}

Tracks the party's public legitimacy.
\begin{itemize}
  \item \textbf{High Mandate:} Allies seek them out, resources flow, doors open.
  \item \textbf{Low Mandate:} Suspicion, bureaucratic obstacles, support withers.
\end{itemize}
Advance the timer when the party gains renown (clean victories, public accolades) or loses face (scandals, broken oaths).

\subsection{Crisis Timer (0--6)}
\label{subsec:crisis}

Tracks the opposition's growing strength.
\begin{itemize}
  \item \textbf{Rising Crisis:} Complications escalate, enemies grow bolder.
  \item \textbf{Managed Crisis:} Breathing room, opportunities to strike back.
\end{itemize}
Advance the timer when the party fails key objectives, ignores looming threats, or the GM spends SB on faction escalation.

\paragraph{Example.}
If the party brokers a peace treaty, advance \textit{Mandate} by 1. If they assassinate a rival leader but leave evidence, advance \textit{Crisis} by 1. When either timer reaches 6, trigger a major faction response (a summons, an assassination attempt, a public trial).

\section{Combat Resource Management}
\label{sec:combat-resources}
\index{combat resources!GM guidance}

In combat, resource management takes on desperate urgency.

\subsection{Supply in Combat}
\label{subsec:combat-supply}

\begin{itemize}
  \item \textbf{Intense Combat:} GM may spend 1 SB to worsen Supply by one step (e.g., Stable $\rightarrow$ Strained).
  \item \textbf{Prolonged Engagement:} Each hour of sustained combat worsens Supply one step.
  \item \textbf{Ammunition Depletion:} Ranged weapons may run low, requiring scavenging actions (e.g., \textbf{Craft} roll to recover arrows).
\end{itemize}

\subsection{Fatigue in Combat}
\label{subsec:combat-fatigue}

\begin{itemize}
  \item Characters apply Fatigue penalties as described in \ref{sec:fatigue-gm}.
  \item Reaching Fatigue level 4+ during combat causes immediate collapse (incapacitated).
  \item Fatigue cannot be cleared during active combat.
\end{itemize}

\subsection{Follower Combat Integration}
\label{subsec:combat-followers}

\begin{itemize}
  \item Followers assist using their Cap (see \ref{subsec:followers-gm}).
  \item Spending \textbf{2+ SB} in combat can worsen a follower's Loyalty or Fitness.
  \item Followers can take combat-relevant independent actions (cost: worsen Fitness one step).
  \item After a follower acts twice in a high-risk combat scene, worsen their Exposure or Fitness.
\end{itemize}

\section{Narrative First: The Fiction Is the Ledger}
\label{sec:resources-narrative-first}
\index{narrative first!resources}

In Fate's Edge, arrows, rations, and waterskins are tracked only in the fiction that surrounds them. Mechanics engage only when those resources become scarce enough to matter. The focus remains always on \textbf{narrative tension} -- the gnawing hunger, the fading light, the last arrow -- not sterile bookkeeping.

\paragraph{Abstracting Coin and Time.}
Routine expenses and petty bribes are handled narratively (\textquotedbl You have enough coin from past jobs to cover a week's stay\textquotedbl). For major expenses (bribing a corrupt official, outfitting an expedition, buying a rare magical component), players may spend \textbf{XP} in place of coin (1 XP = a substantial sum). Use the \textquotedbl XP as currency\textquotedbl\ rule sparingly; it should feel like a meaningful sacrifice, not a tax.

\paragraph{GM Reminder.}
When a player asks \textquotedbl Do I have enough arrows?\textquotedbl\ the answer is almost always \textquotedbl yes\textquotedbl\ until it becomes dramatically interesting to say \textquotedbl no.\textquotedbl\ Similarly, track Supply only when the party is isolated, pursued, or in a hostile environment. Let the world breathe with its own needs and abundances. And when the dice say the world pushes back against mortal plans -- \textbf{listen to what they tell you about the price of ambition.}

\section{Echo Characters: Managing Multi-Character Switching}
\label{sec:echo-characters-gm}
\index{Echo Character}\index{Downtime Log}

Players may choose to temporarily step away from their current character and play another, while their first character continues to act offscreen. This is called an \textbf{Echo Character}. The optional rules in the Player's Guide give the player-facing procedure; this section provides the GM’s mechanical framework for tracking offscreen characters, applying consequences, and maintaining narrative continuity.

\subsection{Core Principles}
\begin{itemize}
  \item Players \textbf{choose} when their character steps offscreen.
  \item Offscreen characters remain active in the fiction and can still affect the world.
  \item Absence creates \textbf{opportunity and risk}, not punishment.
  \item The world reacts to characters even when unseen.
\end{itemize}

\subsection{The Downtime Log}
Each session a character is offscreen, the player completes a simple \textbf{Downtime Log}. The GM records the outcomes and may introduce one consequence per log.

\subsubsection{Log Structure}
The player declares:
\begin{enumerate}
  \item \textbf{Resource Action}: How the character maintains connections or standing (e.g., sends a letter, pays a retainer, visits a contact offscreen).
  \item \textbf{Project Action}: One goal or project advanced (e.g., researches a ritual, trains a follower, builds a contacts list).
  \item \textbf{Event Response}: How the character reacts to unseen events or threats (e.g., avoids a patrol, hides from a rival, investigates a rumor).
\end{enumerate}

\subsubsection{Tier-Based Project Limits}
Players may advance up to a certain number of projects per downtime based on their character’s Tier:
\begin{itemize}
  \item Tier I–II: 1 Project
  \item Tier III–IV: 2 Projects
  \item Tier V: 3 Projects
\end{itemize}

\subsubsection{GM Response}
For each Downtime Log, the GM:
\begin{itemize}
  \item Notes the fictional outcome of each declaration.
  \item May assign \textbf{one meaningful Consequence} (complication, shift, or new opportunity) that will affect the character when they return.
\end{itemize}
Consequences should be specific and story-driven, not punitive. Examples:
\begin{itemize}
  \item \textit{``Your contact has moved; you'll need to find them again.''}
  \item \textit{``The rival you avoided now knows your general location.''}
  \item \textit{``Your project succeeds, but you lost a valuable resource in the process.''}
\end{itemize}

\subsection{Crisis Timer for Offscreen Characters}
Each offscreen character tracks a private \textbf{Crisis Timer [6]}. This timer represents mounting pressure, danger, or narrative trouble that accumulates while the character is away.

\subsubsection{When to Tick the Crisis Clock}
Tick the Crisis Timer +1 when:
\begin{itemize}
  \item The player’s Downtime Log includes a \textbf{failed} or \textbf{partial} outcome (GM discretion).
  \item The GM spends \textbf{2+ Story Beats} on offscreen events related to that character.
  \item The character ignores an obvious threat or obligation (e.g., refuses a Patron’s demand while offscreen).
  \item The character is \textit{In Peril} (status declared by player at start of absence).
\end{itemize}

\subsubsection{When the Crisis Timer Fills}
At 6 segments, the character faces a \textbf{major offscreen crisis}. The GM describes what happens, and the player must address it upon return. Examples:
\begin{itemize}
  \item A rival captures the character’s ally or asset.
  \item The character is wounded or cursed (apply Harm 1 or a Condition).
  \item The character loses a valuable resource (String, Asset tag, or follower Loyalty).
  \item A Patron calls in a debt in an especially inconvenient way.
\end{itemize}
After the crisis, reset the Crisis Timer to 0. The crisis should create a story hook, not remove the character from play.

\subsubsection{Reducing the Crisis Clock}
Players may reduce the Crisis Timer by:
\begin{itemize}
  \item Spending \textbf{1 Boon} (on their active character) to send aid or a message.
  \item Accepting a \textbf{Devil’s Bargain} that introduces a complication on their active character.
  \item Having another player’s active character interact with the offscreen character’s situation (e.g., “I send a scout to check on them”).
\end{itemize}
The GM may also reduce the timer by 1 if the player’s Downtime Log shows proactive risk management.

\subsection{Return Mechanics}
When a player brings an Echo Character back to active play:

\begin{enumerate}
  \item The returning character gains \textbf{2 Boons} for narrative re‑entry (a burst of momentum).
  \item The player rolls \textbf{Wits + Spirit} (DV 3) to reestablish presence. On a success, they reintegrate smoothly. On a failure, the GM applies one additional consequence (e.g., they arrive off-balance, or a complication follows them).
  \item The GM reveals any Consequences accumulated during absence, and the player must address unresolved plot threads.
\end{enumerate}

\subsection{Optional: Buddy System}
If multiple players have offscreen characters, they may nominate a \textbf{Buddy} (another player with an active character) to help track Downtime Logs and Crisis Timers. The Buddy gains \textbf{+1 additional Boon} for the session for this support role (this Boon does not count toward the session cap).

\subsection{Player Guidelines (for GM to communicate)}
Share these with players who wish to use the system:
\begin{itemize}
  \item Declare offscreen goals with clear fictional intent.
  \item Treat absence as character growth, not disappearance.
  \item Expect absence to create tension, opportunity, and drama.
  \item Keep your Downtime Log brief (one sentence per action).
\end{itemize}

\subsection{Full Cross-Reference}
For player-facing rules on Echo Characters (how to declare absence, what a Downtime Log is, and the benefits of returning), see the Player’s Guide, Chapter \ref{ch:assets-followers}. The player version omits Crisis Timers and GM consequence mechanics; those are your tools as GM.